% Simplified plate diagram for the short paper. Topic effects and the swept
% selection parameter are suppressed (stated in the caption).
\begin{tikzpicture}[
    >={Latex[length=1.8mm]}, semithick, font=\small,
    node/.style={circle, thick, minimum size=10mm, align=center,
        inner sep=0.5pt, font=\footnotesize},
    bgN/.style={node, draw=oiGreen, fill=oiGreen!12},
    asN/.style={node, draw=oiOrange, fill=oiOrange!12},
    prN/.style={node, draw=oiBlue, fill=oiBlue!12},
    obs/.style={node, draw=oiGray, fill=black!16},
    plate/.style={rectangle, rounded corners=3pt, draw=black!45, thick,
        inner sep=8pt},
    platelab/.style={font=\scriptsize, text=black, align=left},
    edge/.style={->, draw=black!55},
    lab/.style={font=\scriptsize, text=black},
    newN/.style={node, draw=black!60, fill=black!6},
    hypN/.style={node, draw=black!60, fill=white, dashed, minimum size=7.5mm},
    newedge/.style={->, draw=black!55},
    newlab/.style={font=\scriptsize, text=black},
]

% shared scales (hyperparameters): the top of the hierarchy
\node[hypN] (sb) at (2.7, 9.4) {$\sigma_b$};
\node[hypN] (sv) at (4.8, 9.4) {$\sigma_v$};
\node[hypN] (su) at (6.9, 9.4) {$\sigma_u$};
\node[newlab, align=center] at (0.4, 9.4) {shared\\ scales};

% class plate (partial pooling across the eight subject classes)
\node[newN] (bc) at (2.7, 7.5) {$b_c$};
\node[newN] (vc) at (4.8, 7.5) {$v_c$};
\node[newN] (uc) at (6.9, 7.5) {$u_c$};
\begin{scope}[on background layer]
\node[plate, fit=(bc)(vc)(uc)] (cplate) {};
\end{scope}
\node[newlab, anchor=west] at ($(cplate.east)+(2pt,0)$) {classes $c=1\ldots8$};
\node[newlab, anchor=east, align=center] at ($(cplate.west)-(2pt,0)$)
    {zero-sum};

% quarter-level row
\node[bgN] (g0)    at (1.7, 3.2)  {$g_0(t)$};
\node[asN] (delta) at (3.8, 3.2)  {$\delta_t$};
\node[prN] (pi)    at (5.9, 3.2)  {$\pi_t$};
\node[prN] (rho)   at (8.0, 3.2)  {$\rho_t$};

\node[lab, align=center] at (1.7, 4.15)  {background\\ drift};
\node[lab, align=center] at (3.8, 4.15)  {excess of\\ assisted prose};
\node[lab, align=center] at (5.9, 4.15)  {fraction\\ assisted};
\node[lab, align=center] at (8.0, 4.15)  {chance of\\ declaring};

% paper plate
\node[obs] (L) at (0.6, 0.6)  {$L_i$};
\node[bgN] (mu0) at (2.7, 0.6) {$\mu_{i0}$};
\node[asN] (mu1) at (4.8, 0.6) {$\mu_{i1}$};
\node[prN] (A)  at (6.9, 0.6)  {$A_i$};
\node[obs] (K)  at (4.8, -1.1) {$K_i$};
\node[obs] (D)  at (9.0, 0.6) {$D_i$};
\begin{scope}[on background layer]
\node[plate, fit=(L)(mu0)(mu1)(A)(K)(D)] (pplate) {};
\end{scope}
\node[platelab, anchor=north west] at ($(pplate.south west)+(1pt,20pt)$)
    {papers $i=1\ldots N$\\ $A_i\equiv0$ before 2020};

% dispersions
\node[newN] (ph0) at (3.3, -3.4) {$\phi_0$};
\node[newN] (ph1) at (6.3, -3.4) {$\phi_1$};

% the three background continuations act here
\node[newlab, align=center] (var) at (0, 4.6) {three\\ continuations};
\draw[newedge, dashed] (var) -- (g0);

% edges
\draw[newedge] (sb) -- (bc);
\draw[newedge] (sv) -- (vc);
\draw[newedge] (su) -- (uc);
\draw[edge] (g0)    -- (mu0);
\draw[edge] (L)     -- (mu0);
\draw[edge] (mu0)   -- (mu1);
\draw[edge] (delta) -- (mu1);
\draw[edge] (pi)    -- (A);
\draw[edge] (rho)   -- (D);
\draw[edge] (mu0)   to[out=-55,in=160] (K);
\draw[edge] (mu1)   -- (K);
\draw[edge] (A)     to[out=-125,in=20] (K);
\draw[edge] (A)     -- (D);
\draw[newedge] (bc) -- (mu0);
\draw[newedge] (vc) -- (mu1);
\draw[newedge] (uc) -- (A);
\draw[newedge] (ph0) to[out=30,in=-120] (K);
\draw[newedge] (ph1) to[out=150,in=-60] (K);

% legend, placed clear of the plate label beneath the box
\node[obs, minimum size=5.5mm] at (-4,3.5) {};
\node[anchor=west, font=\scriptsize] at (-3.5,3.5) {observed};
\node[bgN, minimum size=5.5mm] at (-4,2.5) {};
\node[anchor=west, font=\scriptsize] at (-3.5,2.5) {background};
\node[asN, minimum size=5.5mm] at (-4,1.5) {};
\node[anchor=west, font=\scriptsize] at (-3.5,1.5) {assisted};
\node[prN, minimum size=5.5mm] at (-4,0.5) {};
\node[anchor=west, font=\scriptsize] at (-3.5,0.5) {prevalence};
\end{tikzpicture}
